\section{The Company of Heaven and Those We Love}

Scripture never shows a solitary blessed soul. On All Saints the Church reads of ``a great multitude, which no man could number, of all nations and tribes and peoples and tongues, standing before the throne and in sight of the Lamb, clothed with white robes, and palms in their hands.''\endnote{Apoc 7:9, Douay--Rheims (Challoner).} The heavenly city is a bride; the reward is a supper, ``the marriage supper of the Lamb''; and the Epistle to the Hebrews tells the baptized that they have already come ``to mount Sion and to the city of the living God, the heavenly Jerusalem, and to the company of many thousands of angels, and to the church of the firstborn who are written in the heavens, and to God the judge of all, and to the spirits of the just made perfect.''\endnote{Apoc 19:9; 21:2; Heb 12:22--23, Douay--Rheims (Challoner).}

The Church's catechism names heaven by that communion. ``This perfect life with the Most Holy Trinity---this communion of life and love with the Trinity, with the Virgin Mary, the angels and all the blessed---is called `heaven.'\,''\endnote{\emph{Catechism of the Catholic Church} 1024.} And the vision, as the Catechism teaches, is the well-spring of that communion and not its rival. Augustine, in the chapter on the beatific vision, says that God will be so present to the saints ``that we shall see Him by the spirit in ourselves, in one another, in Himself, in the new heavens and the new earth, in every created thing which shall then exist.''\endnote{Augustine, \emph{The City of God} XXII.29, trans.\ Dods (1871).} In one another: the blessed see God in the faces of the blessed.

\subsection{Our dear ones are awaiting us}

The question that comes first to a grieving heart is the simplest. Will I see them again? The Fathers answered it without hesitation. Cyprian of Carthage, writing to his people in a time of plague, when every household had its dead, tells them to lift their eyes from the graves to the homeland:

\begin{quote}
We regard paradise as our country---we already begin to consider the patriarchs as our parents: why do we not hasten and run, that we may behold our country, that we may greet our parents? There a great number of our dear ones is awaiting us, and a dense crowd of parents, brothers, children, is longing for us, already assured of their own safety, and still solicitous for our salvation. To attain to their presence and their embrace, what a gladness both for them and for us in common!\endnote{Cyprian, \emph{On the Mortality} 26, trans.\ Wallis, ANF 5.}
\end{quote}

Cyprian was a bishop and a martyr, and he did not promise his people anything the faith did not hold. The dead in Christ are ``awaiting us''; they are ``longing for us''; and they are ``still solicitous for our salvation,'' caring for those they left. Then he widens the circle to the whole court of heaven: ``There the glorious company of the apostles---there the host of the rejoicing prophets---there the innumerable multitude of martyrs, crowned for the victory of their struggle and passion---there the triumphant virgins \ldots---there are merciful men rewarded, who by feeding and helping the poor have done the works of righteousness.''\endnote{Cyprian, \emph{On the Mortality} 26, trans.\ Wallis, ANF 5.} The family and the saints are one company.

\subsection{Knowing and rejoicing in one another}

Gregory the Great, asked by Peter, the companion of his \emph{Dialogues}, whether ``in heaven the good know the good,'' answers from the parable of the rich man and Lazarus, where Abraham knows Lazarus and the rich man knows both. From that mutual knowledge, he says, ``the reward likewise to both parts increaseth, for the good do more rejoice, when they behold them also in felicity whom before they loved.'' And he adds a wonder:

\begin{quote}
Beside all this, a more wonderful grace is bestowed upon the Saints in heaven: for they know not only them with whom they were acquainted in this world, but also those whom before they never saw, and converse with them in such familiar sort as though in times past they had seen and known one another \ldots\ For seeing they do in that place with unspeakable brightness (common to all) behold God, what is there that they know not, that know him who knoweth all things?\endnote{Gregory the Great, \emph{Dialogues} IV.33, trans.\ P.~W.\ (1608), ed.\ Gardner (1911).}
\end{quote}

The knowledge of friends in heaven is not a distraction from the vision. It is given in the vision, by the light in which everything is known.

Thomas Aquinas, expounding the last article of the Creed, names the fourth thing in which everlasting life consists, after union with God in vision and praise, the fulfilment of desire, and perfect security: it consists \latin{in omnium beatorum iucunda societate}, in the joyful company of all the blessed, which will be supremely delightful because each will have every good together with all the blessed. His reason is charity: \latin{Nam quilibet diliget alium sicut seipsum; et ideo gaudebit de bono alterius sicut de suo. Quo fit ut tantum augeatur laetitia et gaudium unius, quantum est gaudium omnium.}\endnote{Aquinas, \emph{Collationes in Symbolum Apostolorum}, a.~12, Latin text of the Corpus Thomisticum. The English renderings in the text are glosses, not a published translation.} Each will love the other as himself, and so will rejoice in another's good as in his own; and so the joy of each grows by as much as is the joy of all. Alphonsus says the same in his meditation on heaven: ``In that kingdom of love, all love each other tenderly, and each rejoices in the good of the others as if it were his own.''\endnote{Alphonsus, \emph{Preparation for Death}, Consideration XXIX, second point (Duffy, 1843).}

Nor are the bonds of this life abolished. In the \emph{Summa} Thomas asks whether the order of charity remains in heaven, and grants that the blessed will love most those who are nearest to God; but he adds: ``It will however be possible in heaven for a man to love in several ways one who is connected with him, since the causes of virtuous love will not be banished from the mind of the blessed. Yet all these reasons are incomparably surpassed by that which is taken from nighness to God.''\endnote{Aquinas, \emph{Summa theologiae} II-II, q.~26, a.~13, trans.\ English Dominican Province (Project Gutenberg eBook 18755).} The mother will still love her son as her son, and the friend his friend; only now every love is ordered in God and made perfect by him.

\subsection{Their care for us}

The love of the blessed does not stop at the gate of heaven. Paul VI, in the Credo he proclaimed in 1968, confessed that the multitude gathered around Jesus and Mary in paradise forms the Church of heaven, and the Catechism quotes him: there the blessed see God as he is and ``are also, to various degrees, associated with the holy angels in the divine governance exercised by Christ in glory, by interceding for us and helping our weakness by their fraternal concern.''\endnote{\emph{Catechism of the Catholic Church} 1053, quoting Paul VI, \emph{Credo of the People of God} (1968), n.~29.} The Council of Trent teaches it as the Church's faith: the saints, reigning together with Christ, \latin{orationes suas pro hominibus Deo offerre}, offer their prayers to God for men.\endnote{Council of Trent, session XXV (3 December 1563), decree on the invocation of saints, in Denzinger, \emph{Enchiridion symbolorum}, n.~1821 (Latin); likewise the Tridentine Profession of Faith (1564), n.~1867.} Cyprian's ``still solicitous for our salvation'' is the faith of the Church.

Thérèse of Lisieux, in her last illness, made that faith her mission. The epilogue to her story records what she said one evening to Mother Agnes of Jesus: ``I will spend my Heaven in doing good upon earth.'' And she went on at once to say how: ``Nor is this impossible, since from the very heart of the Beatific Vision, the Angels keep watch over us.''\endnote{\emph{Soeur Thérèse of Lisieux}, ed.\ Taylor (1912), Epilogue: ``A Victim of Divine Love.''} The saints do not leave the vision to help us. They help us from inside it.

\subsection{What this asks of us now}

Alphonsus, imagining the moment a faithful soul enters heaven, sets the meeting with family among its first joys: ``What shall be her consolation in rejoining there her relatives or friends, who have been already admitted into Heaven, and in meeting her holy advocates!''\endnote{Alphonsus, \emph{Preparation for Death}, Consideration XXIX, first point (Duffy, 1843).} He tells it as a meditation, in the manner of the saints who teach the faithful to picture what they hope for; the hope itself is the hope of Cyprian and Gregory and Thomas. The Christian who buries a parent, a spouse, a child, may say at the graveside what Cyprian said: they are awaiting us.

Two duties follow, and November keeps both. The first is prayer for those of our dead who may still be in need of purification, the mercy by which the saved are made ready to see; the Church gives the month to it, and \emph{What Death Does Not Finish} sets it out. The second is to love now as the blessed love: in God, and for God, and so for ever. A friendship that is ordered to God here is already part of the company of heaven. The communion we hope to enjoy there is the charity we practise here, made perfect and made permanent.
